\chapter[Estudio cualitativo de funciones]{Estudio cualitativo y representación de funciones}
\label{chap:estudio-funciones}

Representar una función no consiste en encadenar cálculos sin orden. El
objetivo es deducir una descripción global coherente a partir del dominio, los
límites, la continuidad y las derivadas. Cada conclusión debe indicar dónde
vale; en particular, los intervalos separados por puntos excluidos del dominio
no pueden tratarse como uno solo.

\section{Un método de estudio}
\label{sec:metodo-estudio-funciones}

\subsection{Del dominio a la representación cualitativa}
\CanonicalUnit{CM-C05-U001}{}{}

Según la función, no todos los pasos aportan la misma información. Un flujo
razonado de estudio es el siguiente:

\begin{enumerate}
  \item determinar el dominio y sus componentes;
  \item comprobar simetrías o periodicidad;
  \item hallar las intersecciones con los ejes, cuando existan;
  \item estudiar el signo de la función en cada componente relevante;
  \item calcular límites en la frontera del dominio;
  \item estudiar el comportamiento cuando \(x\to+\infty\) y
        \(x\to-\infty\), si procede;
  \item identificar asíntotas y justificar cada una mediante límites;
  \item localizar los puntos donde \(f'\) se anula o no existe;
  \item determinar la monotonía por intervalos;
  \item clasificar los extremos locales y, cuando tenga sentido, los globales;
  \item localizar los puntos donde \(f''\) se anula o no existe;
  \item determinar convexidad y concavidad por intervalos;
  \item confirmar los puntos de inflexión mediante un cambio de convexidad;
  \item reunir la información en tablas de signos o variación coherentes;
  \item bosquejar la gráfica solo después de las deducciones anteriores.
\end{enumerate}

El dominio gobierna todo el análisis. Una ecuación como \(f'(x)>0\) puede
probar crecimiento estricto en cada componente intervalar, pero no permite
comparar valores situados en componentes distintas. Del mismo modo, un punto
donde cambia el signo de una expresión no es un punto de la gráfica si no
pertenece al dominio.

\section{Simetrías y periodicidad}
\label{sec:simetrias-periodicidad}

\subsection{Funciones pares e impares}
\CanonicalUnit{CM-C05-U002}{}{}

Sea \(f:D\to\R\). Para hablar de paridad exigimos que el dominio sea
simétrico respecto de cero:
\[
  x\in D \quad\Longrightarrow\quad -x\in D.
\]
La función es \emph{par} si \(f(-x)=f(x)\) para todo \(x\in D\), e
\emph{impar} si \(f(-x)=-f(x)\) para todo \(x\in D\). La gráfica de una
función par es simétrica respecto del eje vertical; la de una función impar
es simétrica respecto del origen.

Las ecuaciones no son intercambiables. Por ejemplo, \(x^2\) es par y \(x^3\)
es impar. En cambio, la misma expresión \(x^2\) restringida al dominio
\([0,+\infty)\) no se clasifica como par ni impar, porque su dominio no es
simétrico respecto de cero.

\subsection{Funciones periódicas}
\CanonicalUnit{CM-C05-U003}{}{}

Sea \(f:D\to\R\) y sea \(p>0\). Diremos que \(p\) es un \emph{período} de
\(f\) si el dominio es estable por esa traslación,
\[
  D+p=D,
\]
y \(f(x+p)=f(x)\) para todo \(x\in D\). Una función que posee al menos un
período es periódica. Esta definición incluye la compatibilidad del dominio;
la igualdad de valores carecería de sentido si alguno de los argumentos
quedara fuera de él.

Cuando existe un período fundamental, basta estudiar un intervalo que cubra
un ciclo y extender la información por traslación. No toda función periódica
posee un período positivo mínimo: una función constante, por ejemplo, admite
todo \(p>0\).

\section{Intersecciones, límites y asíntotas}
\label{sec:asintotas-estudio}

\subsection{Intersecciones con los ejes}
\CanonicalUnit{CM-C05-U004}{}{}

Las intersecciones con el eje horizontal son los puntos \((x,0)\) para los
que \(x\in D\) y \(f(x)=0\). La intersección con el eje vertical existe solo
si \(0\in D\), y entonces es \((0,f(0))\). Resolver una ecuación sin comprobar
el dominio puede introducir intersecciones inexistentes.

\subsection{Tres tipos de asíntota}
\CanonicalUnit{CM-C05-U005}{}{}

La recta vertical \(x=a\) es una \emph{asíntota vertical} si al menos uno de
los límites laterales \(\lim_{x\to a^-}f(x)\) o
\(\lim_{x\to a^+}f(x)\) es infinito. Debe indicarse el lado y el signo del
límite; no se presupone que ambos coincidan.

La recta horizontal \(y=L\) es una \emph{asíntota horizontal} cuando
\[
  \lim_{x\to+\infty}f(x)=L
  \quad\text{o}\quad
  \lim_{x\to-\infty}f(x)=L.
\]
Los dos infinitos se estudian por separado.

La recta \(y=mx+n\), con \(m\ne0\), es una \emph{asíntota oblicua} en uno
de los dos infinitos si
\[
  \lim_{x\to\pm\infty}\bigl(f(x)-(mx+n)\bigr)=0.
\]
Cuando los límites existen y son finitos, sus coeficientes pueden obtenerse
mediante
\[
  m=\lim_{x\to\pm\infty}\frac{f(x)}{x},
  \qquad
  n=\lim_{x\to\pm\infty}\bigl(f(x)-mx\bigr).
\]
Estas fórmulas son criterios que deben verificarse, no reglas formales
independientes de los límites.

\section{Información de las derivadas}
\label{sec:sintesis-derivadas}

\subsection{Variación, extremos, convexidad e inflexión}
\CanonicalUnit{CM-C05-U006}{}{}

Los criterios de los Capítulos~\ref{chap:derivacion-valor-medio} y
\ref{chap:consecuencias-derivada} se aplican componente a componente:

\begin{itemize}
  \item el signo de \(f'\) determina la monotonía en intervalos;
  \item un cambio de signo de \(f'\) clasifica un extremo local;
  \item los extremos globales exigen comparar todos los candidatos y
        justificar la existencia cuando esta no sea automática;
  \item el signo de \(f''\) determina convexidad o concavidad bajo las
        hipótesis del criterio diferencial;
  \item un punto de inflexión debe pertenecer al dominio y presentar un
        cambio de convexidad.
\end{itemize}

Las ecuaciones \(f'(a)=0\) y \(f''(a)=0\) producen candidatos, no
conclusiones. La primera no garantiza un extremo y la segunda no garantiza un
punto de inflexión. También deben considerarse los puntos pertinentes donde
la derivada correspondiente no existe.

\section{Ejemplo de síntesis}
\label{sec:ejemplo-sintesis-funcion}

\subsection{La función \(f(x)=x-1/x\)}
\CanonicalUnit{CM-C05-U007}{}{}

Consideremos
\[
  f(x)=x-\frac1x.
\]
Su dominio es \(D=\R\setminus\{0\}=(-\infty,0)\cup(0,+\infty)\). Es una
función impar, pues el dominio es simétrico y \(f(-x)=-f(x)\).

Las intersecciones con el eje horizontal se obtienen de
\[
  x-\frac1x=0,
\]
que, dentro del dominio, equivale a \(x^2=1\). Son \((-1,0)\) y \((1,0)\).
No hay intersección con el eje vertical porque \(0\notin D\).

En la frontera del dominio,
\[
  \lim_{x\to0^-}f(x)=+\infty,
  \qquad
  \lim_{x\to0^+}f(x)=-\infty,
\]
de modo que \(x=0\) es una asíntota vertical con comportamientos laterales
opuestos. Además,
\[
  f(x)-x=-\frac1x\longrightarrow0
  \qquad (x\to\pm\infty),
\]
por lo que \(y=x\) es una asíntota oblicua en ambos infinitos.

Las dos primeras derivadas son
\[
  f'(x)=1+\frac1{x^2}>0,
  \qquad
  f''(x)=-\frac2{x^3}.
\]
Por tanto, \(f\) es estrictamente creciente en \(( -\infty,0)\) y en
\((0,+\infty)\). No se afirma que sea creciente en todo \(D\): por ejemplo,
\(-1<1\), pero \(f(-1)=f(1)=0\). No hay puntos críticos ni extremos locales.

Como \(f''(x)>0\) para \(x<0\), la función es convexa en
\(( -\infty,0)\). Como \(f''(x)<0\) para \(x>0\), es cóncava en
\((0,+\infty)\). El cambio se produce al atravesar cero, pero cero no
pertenece al dominio; por ello no existe allí un punto de inflexión.

La representación cualitativa tiene dos ramas crecientes. La rama izquierda
se aproxima a \(y=x\) cuando \(x\to-\infty\), pasa por \((-1,0)\) y tiende a
\(+\infty\) cuando \(x\to0^-\). La rama derecha parte de \(-\infty\) cuando
\(x\to0^+\), pasa por \((1,0)\) y se aproxima a \(y=x\) cuando
\(x\to+\infty\). Esta descripción determina el bosquejo sin usar una gráfica
como sustituto de las deducciones.

\section{Comprobación computacional y alcance}
\label{sec:comprobacion-computacional}

\subsection{Qué puede comprobar una herramienta}
\CanonicalUnit{CM-C05-U008}{}{}

Una calculadora simbólica o una herramienta de representación puede ayudar a
comprobar una derivada, explorar una ventana gráfica o detectar un posible
error de cálculo. No demuestra por sí sola el dominio, la existencia de una
asíntota, la monotonía en un intervalo ni la exhaustividad de los extremos.
La validación matemática debe apoyarse en definiciones, límites y teoremas.

El estudio cualitativo termina aquí con un ejemplo expositivo completo. Los
ejercicios de práctica del material de partida quedan fuera de esta fuente
canónica y de la autorización P4.5B.
